%%%%%%%%%%  scan idx 303  |  folio 290  %%%%%%%%%%
\section*{3. Compléments sur les groupes de Grothendieck des schémas\protect\footnote{Les résultats de ce numéro ne seront pas utilisés dans les Exp.\ V à VIII.}}

\textbf{Proposition 3.1.0.} Soit
\[
\begin{tikzcd}
X \arrow[d,"f"'] & X' \arrow[l,"g"']\arrow[d,"f'"]\\
Y & Y' \arrow[l,"h"']
\end{tikzcd}
\]
un carré cartésien de schémas, où \(Y\) est quasi-compact et \(f\) quasi-compact et quasi-séparé, et soit \(E\in\Ob D^b(X)\) un complexe à cohomologie quasi-cohérente. On suppose \(X\) et \(Y'\) tor-indépendants sur \(Y\) (III 1.5). On suppose d'autre part qu'on a ou bien \(h\) de tor-dimension finie ou bien \(E\) d'amplitude plate finie rel. à \(f\). Dans ces conditions, la flèche de changement de base (cf. SGA 4 XVII 4)
\[
 Lh^*\,Rf_*E\longrightarrow Rf'_*\,Lg^*E
\]
est définie et est un isomorphisme.

\emph{Preuve}~: Pour définir la flèche, on utilise la formule de dualité triviale (SGA 4 XVII 2.3) pour \(h\), qui dit que l'on a un isomorphisme fonctoriel canonique
\begin{equation}\tag{$*$}
 \Hom(Lh^*L,M)\xrightarrow{\sim}\Hom(L,Rh_*M)
\end{equation}
pour \(L\in\Ob D^-(Y)\), \(M\in\Ob D^+(Y')\). Utilisant (III 3.5.2) et (III 3.7.1), on vérifié facilement qu'on a \(Lg^*E\in\Ob D^b(X')\) et \(Rf_*E\in\Ob D^-(Y)\), de sorte qu'appliquant (*) avec \(L=Rf_*E\) et \(M=Rf'_*Lg^*E\), on est ramené à définir une flèche
\[
 Rf_*E\longrightarrow Rf_*Rg_*Lg^*E.
\]
On prend pour celle-ci l'image par \(Rf_*\) de la flèche d'adjonction
\[
 E\longrightarrow Rg_*Lg^*E,
\]
provenant de la formule de dualité triviale pour \(g\). La flèche de changement de base 
